% ======================================================================
% Core Vision Perfect V1 — Báo cáo giải pháp sơ tuyển AIC 2026
% Biên dịch: pdflatex -> bibtex main -> pdflatex -> pdflatex  (xem report/README.md)
% Template: ai_conquer2026.cls (AIVN). KHÔNG sửa cls/sty — mọi tùy biến đặt ở đây.
% ======================================================================
\documentclass{ai_conquer2026}

\addbibresource{references.bib}

% ---- Logo đội (tùy chọn) --------------------------------------------
% cls mặc định tìm Figures/logo.pdf. Ta override an toàn: nếu chưa có
% figures/logo.pdf thì bỏ khối logo để tài liệu vẫn biên dịch sạch.
\pretitle{%
  \begin{flushleft}
    \IfFileExists{figures/logo.pdf}{\includegraphics[height=2.0cm]{figures/logo.pdf}\par}{}%
  \end{flushleft}
  \begin{center}\LARGE\bfseries
}
\posttitle{\end{center}}

\title{Core Vision Perfect V1 --- Trợ lý truy xuất khoảnh khắc video tiếng Việt cho AIC 2026}
\author{%
  % TODO(đội): tên đội + danh sách thành viên + đơn vị, theo yêu cầu của BTC.
  Đội \texttt{TODO} \\ \normalsize TODO Thành viên 1 · TODO Thành viên 2 · TODO Thành viên 3
}
\date{Vòng sơ tuyển --- Hội thi Thử thách Trí tuệ Nhân tạo TP.HCM 2026 \\ \normalsize \today}

\begin{document}
\selectlanguage{vietnamese}
\maketitle

\begin{minisummaryboxaivn}
\textbf{Tóm tắt.} Core Vision Perfect V1 là hệ thống trợ lý truy xuất khoảnh khắc
video quy mô lớn (hàng trăm đến $\sim$1000+ giờ) bằng truy vấn tiếng Việt, phục vụ
cả hai hình thức thi 2026: \emph{tương tác} (người + trợ lý) và \emph{tự động}
(trợ lý đấu trợ lý). Lõi hệ thống là ensemble hai lane embedding
(SigLIP-2 đa ngữ + PE-Core tiếng Anh)~\cite{tschannen2025siglip2,bolya2025pe},
hợp nhất với BM25 bốn kênh (OCR/ASR/caption/metadata), tinh chỉnh SuperGlobal~\cite{shao2023superglobal},
rerank cross-encoder tùy chọn (BLIP-2 ITM~\cite{li2023blip2} / Qwen3-VL-Reranker),
quy hoạch động DANTE cho TRAKE, MMR~\cite{carbonell1998mmr} cho AVS, và VQA đa khung hình
cho Q\&A. Toàn bộ pipeline offline có thể resume, online trả lời
dưới 1 giây/truy vấn; 596 test CPU + CI đảm bảo tính đúng của công thức chấm
và định dạng nộp bài. % TODO: cập nhật kết quả sơ tuyển khi có.
\end{minisummaryboxaivn}

\tableofcontents
\newpage

% ======================================================================
\section{Giới thiệu và bài toán}
% ======================================================================

Hội thi Thử thách Trí tuệ Nhân tạo TP.HCM 2026 (AIC 2026) đặt một bài toán duy nhất:
xây dựng \emph{trợ lý ảo thông minh} truy xuất thông tin trên kho dữ liệu đa phương tiện
lớn --- tin tức truyền hình và, theo định hướng chủ đề 2026, có thể bổ sung video
sousveillance/egocentric --- với quy mô có thể vượt 1000 giờ. Điểm mới quan trọng nhất
của năm 2026 là hai \emph{hình thức thi}: (i) \textbf{tương tác} --- thể thức truyền thống,
con người thao tác trên công cụ của đội, trong đó việc dùng agent tự động bên trong công cụ
là \emph{tùy chọn}; và (ii) \textbf{tự động} (thí điểm) --- trợ lý đấu trợ lý, không người
can thiệp; BTC \emph{chưa công bố} đặc tả kỹ thuật của hình thức này, nhưng chắc chắn
không đơn thuần là ``tự động nộp file''. Hệ thống của chúng tôi vì vậy được thiết kế
\emph{một lõi hai vỏ}: cùng một engine truy xuất phục vụ ứng dụng tương tác (Streamlit)
lẫn dịch vụ HTTP/JSON máy-gọi-máy (mục~\ref{sec:autoagent}) sẵn sàng thích ứng khi spec
hình thức tự động được công bố.

Về thời gian: BTC phát hành dữ liệu, đề bài, baseline và độ đo \emph{chậm nhất 25/07/2026};
sơ tuyển diễn ra tháng 8/2026 trên Codabench (kết quả 30/8); chung kết on-site 12--26/9/2026;
trao giải tháng 10/2026. Vòng sơ tuyển yêu cầu nộp kèm \emph{báo cáo giải pháp} --- chính là
tài liệu này. Các dạng đề: KIS (mô tả văn bản $\rightarrow$ đúng một khoảnh khắc),
KIS-V (xem clip ngắn chiếu trên màn hình --- chỉ được \emph{xem}, cấm quay chụp, âm thanh
có thể bị tắt), và Q\&A là chắc chắn; TRAKE (chuỗi sự kiện có thứ tự trong một video)
có tiền lệ 2025; AVS \emph{chưa chắc} xuất hiện nên được giữ sau cờ cấu hình; KIS-C
(hội thoại tiết lộ dần) hiện là \emph{phong cách hệ thống} được khuyến khích, chưa phải
dạng đề chính thức. FAQ của BTC xác nhận \emph{không giới hạn} mô hình, thuật toán
hay công cụ.

Định dạng nộp sơ tuyển (theo đặc tả 2025, cổng thi xác nhận 2026 giữ nguyên định dạng):
mỗi truy vấn một file CSV UTF-8 \emph{không header}, tối đa 100 dòng, phân tách bằng dấu phẩy;
KIS: \texttt{video\_name,frame\_id}; Q\&A: \texttt{video\_name,frame\_id,answer} (answer
$\le$100 ký tự, tiếng Việt hoặc tiếng Anh, bọc nháy kép nếu chứa dấu phẩy); TRAKE:
\texttt{video\_name,f1,\ldots,fN}. Gói nộp là file zip \emph{bắt buộc} chứa thư mục tên
\texttt{submission}. Điểm mỗi truy vấn là trung bình của $\max$ R-Score trong top-$k$ với
$k\in\{1,5,20,50,100\}$ --- kết quả đúng càng đứng đầu bảng xếp hạng điểm càng cao. Mỗi đội
có tối đa 20 lượt nộp, không quá 5 lượt/ngày (năm 2025 còn giới hạn khung giờ nộp buổi sáng
9:00--11:59; chúng tôi vận hành theo kịch bản chặt hơn để an toàn). Ở chung kết, truy vấn là
văn bản hoặc clip chỉ-được-xem, chấm theo server kiểu DRES~\cite{rossetto2021dres}
(bản hiện hành 2.0.4): nộp đúng càng sớm điểm càng cao, nộp sai bị trừ điểm; năm 2025 KIS có
5 phút với 5 gợi ý nhỏ giọt mỗi phút, KIS-V có 4 phút với clip 20 giây.

% ======================================================================
\section{Kiến trúc hệ thống}
% ======================================================================

Chúng tôi quy bài toán về \emph{truy xuất ảnh mức keyframe}: BTC đảm bảo mọi đoạn đáp án
chứa ít nhất một keyframe, nên tìm đúng keyframe là ăn trọn điểm. Kiến trúc tách làm hai
pha. \textbf{Pha offline} (build một lần, mọi bước resumable): video qua TransNetV2 (fallback
PySceneDetect) thành keyframes; keyframes được mã hóa song song bởi các lane embedding
(mục~\ref{sec:lanes}) vào FAISS; đồng thời trích OCR (EasyOCR vi+en), ASR (PhoWhisper),
caption tiếng Việt (Vintern-1B) và metadata thành chỉ mục BM25 \emph{persisted}; nhãn object
(Open Images) được nén thành parquet compact. \textbf{Pha online} (mỗi truy vấn dưới 1 giây):
truy vấn tiếng Việt được Gemini dịch + mô tả thị giác + sinh biến thể (degrade tự động
Gemini $\rightarrow$ Google Translate $\rightarrow$ raw khi mất mạng); các lane encode mọi
biến thể, lấy top-$K$ FAISS, max-fuse theo biến thể, tinh chỉnh SuperGlobal, hợp nhất lane
theo trọng số, cộng tín hiệu BM25 + object boost + neighbor boost trên \emph{đúng tập ứng
viên} (độ phức tạp $O(K)$), rồi tùy chọn rerank cross-encoder trước khi xuất kết quả/CSV/DRES.

\begin{figure}[H]
  \centering
  \fbox{\parbox{0.92\linewidth}{\centering\vspace{3em}
    \textbf{TODO HÌNH 1:} Sơ đồ pipeline offline/online hai cột\\
    (xuất từ \texttt{docs/PROJECT\_CONTEXT.md} mục 2 sang draw.io/TikZ,\\
    lưu \texttt{figures/pipeline.pdf} rồi thay khối này bằng \texttt{\textbackslash includegraphics})
    \vspace{3em}}}
  \begin{figurecaptionmodern}
    \capfig{Kiến trúc hai pha của Core Vision Perfect V1. \textbf{TODO: chèn hình.}}
  \end{figurecaptionmodern}
\end{figure}

Bất biến trung tâm là \emph{catalog}: \texttt{global\_id} $=$ số dòng manifest.parquet $=$
số dòng FAISS index; index lưu chữ ký corpus và \texttt{model\_tag}, lệch là từ chối phục vụ
--- vì sai một hàng là sai \texttt{frame\_idx}, tức 0 điểm. \texttt{frame\_idx} nộp bài là số
frame trong video \emph{gốc} lấy từ \texttt{map-keyframes}; Batch~1 sơ tuyển 2026 phát
hành ĐỦ 873/873 map csv chính thức (đã kiểm kê từng file), và với các gói tương lai
chỉ-có-video, công cụ
\texttt{scripts/05\_rebuild\_map\_keyframes.py} tái tạo chúng bằng dhash từng keyframe, quét
video theo bước nhảy rồi khớp đơn điệu bằng quy hoạch động (thứ hạng keyframe có thứ tự thời
gian nên ràng buộc được các shot lặp), cuối cùng tinh chỉnh chính xác trong cửa sổ
$\pm$stride. Toàn bộ cấu hình nằm trong \texttt{configs/settings.yaml}, override qua biến
môi trường tiền tố \texttt{CVP\_} (ví dụ \texttt{CVP\_SEARCH\_\_RERANKER=qwen\_reranker}) ---
mọi thí nghiệm trong mục~\ref{sec:exp} đều là thay đổi cấu hình thuần, không sửa mã.

% ======================================================================
\section{Các thành phần chính}
% ======================================================================

\subsection{Ensemble hai lane embedding (+ lane tùy chọn)}\label{sec:lanes}
Lane mặc định gồm \texttt{google/siglip2-so400m-patch16-384}~\cite{tschannen2025siglip2}
(đa ngữ, đọc tiếng Việt trực tiếp, đồng thời là lane duy nhất được fine-tune ---
mục~\ref{sec:training}) và \texttt{PE-Core-bigG-14-448}~\cite{bolya2025pe} qua open\_clip
(lane tiếng Anh mạnh nhất 2026, đúng model đội 79/88 điểm AIC 2025 sử dụng; fallback tự động
PE-Core-L $\rightarrow$ DFN5B khi thiếu VRAM). Điểm hai lane được min-max rồi cộng trọng số
0.55/0.45 --- bộ trọng số đã chứng minh, có thể chuyển \texttt{search.fusion\_method: rrf}
khi chưa kịp tune~\cite{bruch2023fusion}. Ba lane \emph{tùy chọn} phục vụ ablation và dự
phòng: \texttt{Qwen/Qwen3-VL-Embedding-2B} (tiếng Việt bản địa, cắt MRL 1024-d),
\texttt{jinaai/jina-clip-v2} (89 ngôn ngữ, Matryoshka) và
\texttt{facebook/metaclip-2-worldwide-huge-quickgelu} (SOTA đa ngữ XM3600) --- bật bằng cách
thêm tên lane vào \texttt{embedding.ensemble\_members}. Mỗi index ghi \texttt{model\_tag}
để phát hiện lệch checkpoint giữa máy build và máy thi.

\subsection{BM25 bốn kênh, object boost và SuperGlobal}
Chỉ mục văn bản BM25+ được tokenize \emph{một lần} lúc build và persist
(\texttt{artifacts/text\_index/*.json.gz}); lúc truy vấn chỉ chấm điểm trên tập ứng viên dense
nên chi phí là $O(K\cdot|q|)$ thật sự. Bốn kênh --- OCR, ASR, caption, metadata --- cùng
object boost (từ vựng 228 danh từ tiếng Việt ánh xạ 600 lớp Open Images, parser đa ràng buộc
kiểu ``2 người bên trái và 1 người bên phải'') được hợp nhất bằng weighted-sum theo
\texttt{search.weights}; chỉ mục object dạng parquet compact (build bằng
\texttt{scripts/03\_build\_aux\_indexes.py --objects-index}) được ObjectBooster tự ưu tiên khi
có mặt. SuperGlobal~\cite{shao2023superglobal} tinh chỉnh top-$K$ \emph{trên mọi biến thể truy
vấn} (max-fuse, mask hàng zero) --- biến thể mở rộng tìm ra hit không còn bị rerank dìm.

\subsection{TRAKE bằng DANTE và AVS bằng MMR}
Với TRAKE, từng sự kiện được encode bởi \emph{mọi} lane, pool $\le$30 video, rồi quy hoạch
động DANTE $O(N\cdot T)$: $DP[j,t] = sim[t,j] + \max_{\tau<t}\bigl(DP[j{-}1,\tau] -
\lambda\,(time_t - time_\tau)\bigr)$ với running-max trượt, ràng buộc khoảng cách
$[min\_gap, max\_gap]$ và $\lambda = 0.0005$/giây --- đúng thuật toán đạt ``Outstanding
TRAKE'' AIC 2025; \texttt{temporal.algo: beam} giữ phương án beam đa dạng để A/B. Parser đề
TRAKE đọc đúng định dạng đề thật của BTC (dòng ngữ cảnh + sự kiện tiền tố \texttt{E1:}\ldots).
Với AVS (giữ sau cờ vì 2026 chưa chắc thi), kết quả được đa dạng hóa bằng
MMR~\cite{carbonell1998mmr} mức embedding ($\lambda=0.7$) khử near-duplicate \emph{xuyên}
video --- tin tức tái sử dụng b-roll rất nhiều --- kèm trần 3 dòng/video và giãn cách 10 giây.

\subsection{Rerank cross-encoder pairwise}
Tầng rerank mới của Perfect V1: \texttt{search.reranker} nhận \texttt{blip2\_itm}
(đầu ITM của BLIP-2~\cite{li2023blip2}, đúng blueprint đội 76.4/88 AIC 2025) hoặc
\texttt{qwen\_reranker} (\texttt{Qwen/Qwen3-VL-Reranker-2B}, model 01/2026, chấm cặp
ảnh--văn bản đa ngữ trực tiếp bằng tiếng Việt, không cần dịch). Reranker chỉ chấm lại
top đầu bảng, lỗi model thì giữ nguyên thứ hạng --- không bao giờ làm kết quả tệ hơn vì
sự cố. Tầng này độc lập với VLM listwise rerank (\texttt{search.vlm\_rerank}, Gemini chấm
0--10) đã có từ bản kế thừa; hai cơ chế được so trực tiếp trong ablation A9.

\subsection{VQA đa khung hình cho Q\&A}
Ứng viên Q\&A được gộp nhóm theo (video, ngắt cảnh), VQA trả lời \emph{theo từng nhóm} để
mỗi dòng CSV mang đúng answer của nhóm nó --- yêu cầu bắt buộc của công thức chấm. Nâng cấp
2026: \texttt{vqa.frames\_per\_answer} (mặc định 3) gửi \emph{dải} khung hình lân cận trong
một lần gọi kèm bỏ phiếu nhất quán, khắc phục lớp câu hỏi ``giải toán trong video'' của đề
2025 mà một khung hình đơn không đủ ngữ cảnh. Chuỗi model VQA/enhance dùng Gemini với
fallback tự động \texttt{gemini-3.5-flash} $\rightarrow$ \texttt{gemini-3-flash-preview}
$\rightarrow$ \texttt{gemini-2.5-flash}, offline thì rơi về Vintern-1B; answer luôn được
sanitize về $\le$100 ký tự đúng thể lệ.

\subsection{Tăng cường truy vấn thời gian và retry tự tin thấp}
Hai knob phục vụ đề khó: \texttt{search.temporal\_boost} tách truy vấn dạng
``\ldots{} sau khi \ldots'' thành Before/Now/After và thưởng điểm các frame có ngữ cảnh
láng giềng khớp (ý tưởng Vortex, đội 79.6/88 AIC 2025); \texttt{search.low\_confidence\_retry}
--- dành cho hình thức tự động --- phát hiện bảng xếp hạng ``phẳng'' (margin thấp), tự
reformulate truy vấn và trộn RRF hai lượt tìm.

\subsection{Auto-agent và dịch vụ máy-gọi-máy}\label{sec:autoagent}
Cho hình thức tự động, \texttt{pipeline/auto\_agent.py} chạy trọn gói: phân loại dạng đề,
tìm kiếm, gate tự tin theo margin, mở rộng thích ứng, reformulate retry, STAR-style
đa bước cho Q\&A, log tương tác episodic; adapter giao thức được trừu tượng hóa chờ spec
BTC. Nền tảng máy-gọi-máy là dịch vụ FastAPI khởi động bằng \texttt{cvp serve}
(console entry \texttt{cvp} khai báo trong \texttt{pyproject.toml}) với các endpoint
\texttt{/health}, \texttt{/search/text}, \texttt{/search/qa}, \texttt{/search/trake},
\texttt{/search/avs}, \texttt{/nearest/\{global\_id\}}, \texttt{/keyframe/\{global\_id\}}.
Ứng dụng tương tác Streamlit giữ 5 tab nghiệp vụ, đồng hồ 5 phút KIS-C realtime, basket,
export CSV validate sẵn và feedback Rocchio trong không gian từng lane.

% ======================================================================
\section{Huấn luyện bổ trợ: LoRA-LiT + WiSE-FT (kênh phụ)}\label{sec:training}
% ======================================================================

Bằng chứng từ AIC/VBS nhiều năm cho thấy các đội thắng đều dùng \emph{zero-shot ensemble +
dịch câu}; vì vậy fine-tune trong hệ thống này là \emph{kênh bổ sung có đo đạc}, không phải
mặc định. Dữ liệu huấn luyện gồm caption tiếng Việt do Vintern-1B sinh cho chính keyframe
của corpus (in-domain, giá trị nhất) cộng hai bộ công khai KTVIC (21.6k) và UIT-ViIC (19.3k)
qua \texttt{scripts/12\_build\_public\_datasets.py}; lọc độ dài 15--400 ký tự, khử near-dup
($\cos>0.97$), và \emph{split theo video} chống rò rỉ. Cơ chế LiT: đóng băng image tower
(FAISS index giữ nguyên --- không phải re-embed), LoRA~\cite{hu2022lora} $r{=}32$ trên text
tower của SigLIP-2, giữ logit\_scale/bias pretrained; loss sigmoid SigLIP mặc định
(knob \texttt{loss: infonce} để A/B), distillation anchor $1-\cos$(student, teacher) với
$\beta: 0.3\rightarrow0$, trộn 25\% cặp anchor tiếng Anh chống quên đa ngữ, và tùy chọn hard
negative cùng video. Sau huấn luyện, WiSE-FT~\cite{wortsman2022wiseft} quét
$\alpha\in\{0.4,0.5,0.6,1.0\}$ nội suy base$\leftrightarrow$tuned trên val R@5 ---
$\alpha{=}1.0$ chính là tower gốc nên winner không bao giờ tệ hơn zero-shot. Trên Colab
H100, autopilot dò micro-batch theo OOM thực nghiệm (effective batch $\approx$2048, bf16),
checkpoint atomic lên Drive và \emph{resume chính xác giữa epoch}. Kênh này chỉ được bật thi
đấu (\texttt{CVP\_EMBEDDING\_\_MODEL=finetuned} hoặc thêm vào ensemble) sau khi
\texttt{scripts/eval\_model.py} xác nhận thắng baseline trên bộ đề dev --- kết quả đo tại
mục~\ref{sec:exp} (hàng A5). % TODO: điền kết quả train khi notebook 02 chạy xong với data 2026.

% ======================================================================
\section{Thực nghiệm}\label{sec:exp}
% ======================================================================

Bộ đề phát triển hiện hành là gói 89 truy vấn chung kết AIC 2025 nguyên bản (73 KIS, 9 Q\&A,
7 TRAKE) chạy trên corpus cục bộ; khi BTC phát hành dữ liệu và baseline 2026 (chậm nhất
25/07), toàn bộ bảng dưới sẽ được đo lại bằng \emph{một lệnh} trên scorer chính thức
(\texttt{cvp/eval/official.py} --- cài đúng công thức R@k của BTC, 60+ test):

\begin{aivnoutputbox}[Tái lập toàn bộ bảng ablation]
python scripts/26_run_ablations.py --query-dir queries/dev-2025-finals \
    --gt queries/dev-2025-finals/gt.json  # gt.json: đội tự soạn cho câu đã xác minh
# chạy chọn lọc: --only A1 A9 A10 ; kết quả lưu artifacts/ablations/
\end{aivnoutputbox}

\begin{table}[H]
\centering
\begin{adjustbox}{max width=\linewidth}
\begin{tabular}{@{}llcccc@{}}
\toprule
\textbf{ID} & \textbf{Cấu hình (nhãn trong scripts/26)} & \textbf{R@1} & \textbf{R@5} & \textbf{R@100} & \textbf{Final} \\
\midrule
A1 & siglip2-only vs ensemble(siglip2+openclip) & TODO & TODO & TODO & TODO \\
A2 & SuperGlobal off/on & TODO & TODO & TODO & TODO \\
A3 & visual-only / all-signals / rrf-fusion & TODO & TODO & TODO & TODO \\
A4 & TRAKE: beam / dante / dante-ensemble & TODO & TODO & TODO & TODO \\
A5 & zero-shot vs LoRA-LiT+WiSE-FT (notebook 02) & TODO & TODO & TODO & TODO \\
A6 & VLM listwise rerank off/gemini & TODO & TODO & TODO & TODO \\
A7 & AVS: greedy / MMR $\lambda{=}0.7$ / $\lambda{=}0.5$ & TODO & TODO & TODO & TODO \\
A8 & QA: 1 answer toàn cục vs per-group & TODO & TODO & TODO & TODO \\
A9 & cross-rerank: off / blip2 / qwen & TODO & TODO & TODO & TODO \\
A10 & temporal-boost off/on + low-conf-retry & TODO & TODO & TODO & TODO \\
\bottomrule
\end{tabular}
\end{adjustbox}
\captable{Ablation A1--A10 trên bộ đề dev 89 truy vấn. \textbf{TODO: điền số từ
\texttt{artifacts/ablations/} sau khi chạy \texttt{scripts/26\_run\_ablations.py} trên
artifacts thật.} A5 in bảng riêng trong notebook 02; quét trọng số đầy đủ của A3 nằm ở
\texttt{scripts/21\_tune\_weights.py}.}
\end{table}

Trọng số tín hiệu được tune \emph{offline} trên score map đã cache: \texttt{scripts/23\_dump\_signals.py}
chạy engine một lần mỗi câu, \texttt{scripts/21\_tune\_weights.py} tối ưu rồi in đúng các dòng
\texttt{CVP\_SEARCH\_\_WEIGHTS\_\_*} để dán sang máy thi --- không tốn lượt nộp nào của Codabench.
Độ trễ được đo bằng \texttt{scripts/50\_bench\_latency.py} (mặc định $N{=}200$ truy vấn, báo cáo
p50/p95/p99 của đường nóng \texttt{search\_text}); chỉ tiêu vận hành đã chốt: \textbf{p50
$\le$ 200 ms, p95 $\le$ 500 ms} trên laptop thi đấu --- dư ngân sách lớn trong khung 4--5 phút
mỗi truy vấn chung kết. % TODO: dán bảng percentile thật sau khi build artifacts 2026.
Chất lượng kỹ thuật của mã được găm bằng 596 test CPU thuần (không cần GPU/mạng/data thật)
chạy trong CI, phủ scorer chính thức, packager, parser đề thật của BTC và toàn bộ engine
tích hợp với fake model.

% ======================================================================
\section{Vận hành thi đấu và bài học}
% ======================================================================

Quy trình vòng loại là một lệnh duy nhất, thiết kế để \emph{tiết kiệm lượt nộp} (20 tổng,
5/ngày --- và có thể bị bó vào khung giờ sáng như 2025):

\begin{aivnoutputbox}[Vòng loại: chạy đề $\rightarrow$ validate $\rightarrow$ zip $\rightarrow$ chấm thử offline]
python scripts/20_run_queries.py --query-dir <thư mục đề> --zip --gt <gt.json nếu có>
# zip đầu ra CHỨA folder 'submission' đúng thể lệ Codabench, kèm MANIFEST sha256;
# --gt in bảng điểm per-query/per-task/MEAN FINAL đúng công thức BTC trước khi nộp.
\end{aivnoutputbox}

Packager validate mọi CSV theo thể lệ (không header, $\le$100 dòng, answer $\le$100 ký tự,
quote khi có dấu phẩy) \emph{trước} khi đóng gói --- lỗi định dạng bị chặn tại chỗ thay vì
tốn một lượt nộp. Mọi lượt nộp thật đều đi sau ít nhất một lần chấm thử offline bằng
\texttt{--gt} trên bộ dev; kỷ luật nộp: không quá 2 lượt cho ngày đầu (dò phân phối đề),
giữ tối thiểu 3 lượt cho ngày cuối. Ở chung kết, laptop chạy app + artifacts đồng bộ về từ
Drive; DRES client (\texttt{submission/dres\_client.py}, urllib thuần, không bao giờ tự retry
lệnh bị server từ chối) trỏ \texttt{submission.dres\_base\_url} khi BTC công bố endpoint;
đồng hồ 5 phút với nhịp gợi ý 1 phút/lần đã dựng sẵn trong app theo đúng thể thức KIS 2025.
Trước mỗi ngày thi chạy lại \texttt{scripts/doctor.py} (kiểm artifacts/chữ ký/model) và
\texttt{scripts/50\_bench\_latency.py}. Bài học lớn nhất từ các vòng review đối kháng nội bộ:
(1) sai \texttt{frame\_idx} do thiếu map-keyframes là lỗi \emph{chết người} --- phải có công
cụ tái tạo (mục 2) thay vì tin dữ liệu tải về; (2) parser đề thật quan trọng ngang model ---
một dòng ngữ cảnh TRAKE đọc sai làm lệch cả chuỗi frame về 0 điểm; (3) mọi tầng thông minh
(Gemini, reranker, VLM) đều phải có đường degrade để hệ thống không bao giờ chết giữa trận.
% TODO: bổ sung bài học thực chiến sau sơ tuyển 8/2026.

% ======================================================================
\section{Kết luận}
% ======================================================================

Core Vision Perfect V1 hợp nhất các cơ chế đã được kiểm chứng bởi những đội dẫn đầu AIC 2025
--- ensemble PE-Core, RRF/Rocchio, BLIP-2 ITM rerank, DANTE TRAKE, dense keyframes ---
với các nâng cấp 2026: rerank pairwise Qwen3-VL đọc thẳng tiếng Việt, VQA dải khung hình,
temporal boost, retry tự tin thấp, và một lớp dịch vụ máy-gọi-máy đón đầu hình thức thi tự
động. Triết lý xuyên suốt là \emph{đúng trước, nhanh sau, thông minh cuối}: bất biến catalog
và scorer chính thức găm tính đúng; chỉ mục persisted và ngân sách trễ p95 $\le$ 500 ms găm
tốc độ; các tầng LLM/VLM chỉ xếp lại đầu bảng và luôn có đường lui. Toàn bộ thí nghiệm là
thay đổi cấu hình thuần có thể tái lập bằng một lệnh, nên khi dữ liệu 2026 phát hành
(chậm nhất 25/07) hệ thống chỉ cần build lại artifacts và điền số vào các bảng TODO của báo
cáo này. Hướng tiếp theo: hoàn thiện adapter cho spec hình thức tự động ngay khi BTC công bố,
đo lại ablation trên dữ liệu 2026 (đặc biệt nếu xuất hiện lifelog/egocentric), và nộp phiên
bản mở rộng của báo cáo cho SOICT 2026. % TODO: cập nhật kết quả sơ tuyển thực tế.

\newpage
\section*{Tài liệu tham khảo}
\printbibliography[heading=none]

\end{document}
